\documentclass[11pt,a4paper]{article}
\usepackage[utf8]{inputenc}
\usepackage[vietnamese]{babel}
\usepackage{amsmath,amssymb}
\usepackage[version=4]{mhchem}
\usepackage{geometry}
\geometry{top=0.9cm,bottom=0.9cm,left=1.35cm,right=1.35cm,headheight=24pt,headsep=5pt,footskip=11pt}
\usepackage{graphicx}
\usepackage{xcolor}
\usepackage{tabularx}
\usepackage{array}
\usepackage{tikz}
\usepackage{multicol}
\usepackage{fancyhdr}
\usepackage{enumitem}

\definecolor{hfRed}{HTML}{D32F2F}
\definecolor{hfGreen}{HTML}{8BC390}
\definecolor{mainGreen}{HTML}{2E7D32}

\usepackage{ex_test}

\pagestyle{fancy}
\fancyhf{}
\renewcommand{\headrulewidth}{0.75pt}
\renewcommand{\headrule}{\hbox to\headwidth{\color{hfGreen}\leaders\hrule height \headrulewidth\hfill}}
\renewcommand{\footrulewidth}{0.75pt}
\renewcommand{\footrule}{\hbox to\headwidth{\color{hfGreen}\leaders\hrule height \footrulewidth\hfill}}

\fancyhead[L]{\small\itshape\color{hfRed}Tài liệu lưu hành nội bộ}
\fancyhead[R]{\small\color{hfRed}Thầy \textbf{TRẦN VĂN THẠNH} - 0777.470.803}

\fancyfoot[L]{\footnotesize\itshape\color{hfRed}CS1: 6/15 Nguyễn Hoàng, P. Kim Long, TP Huế \quad|\quad CS2: 24 Đặng Thái Thân, TP Huế}
\fancyfoot[R]{\small\textbf{\color{hfRed}Trang \thepage}}

\setlength{\parindent}{0pt}
\setlength{\parskip}{0.0pt}

\def\loigiai#1{}

\begin{document}
\Opensolutionfile{ans}[ans-de15-hs]

\noindent
\begin{minipage}[t]{0.48\textwidth}
	\textbf{SỞ GIÁO DỤC VÀ ĐÀO TẠO THỪA THIÊN HUẾ}\\
	\textbf{TRƯỜNG THPT HAI BÀ TRƯNG}\\
	\textbf{ĐỀ CHÍNH THỨC}\\
	\textit{(Đề thi có 04 trang)}
\end{minipage}\hfill
\begin{minipage}[t]{0.48\textwidth}
	\raggedleft
	\textbf{\color{mainGreen}ĐỀ KIỂM TRA ĐỊNH KỲ CHƯƠNG 5 (PIN ĐIỆN – ĐIỆN PHÂN)}\\
	\textbf{MÔN: HOÁ HỌC -- LỚP 12}\\
	\textit{Thời gian làm bài: 50 phút (không kể thời gian phát đề)}\\
	\textbf{Mã đề thi: 115}
\end{minipage}


\vspace{1pt}
\noindent\textit{Họ, tên thí sinh: \dots\dots\dots\dots\dots\dots\dots\dots\dots\dots\dots\dots\dots\dots\dots\dots\dots\dots\dots\dots\dots\dots Lớp: \dots\dots\dots\dots SBD: \dots\dots\dots\dots}\\
\noindent\textit{(Cho biết nguyên tử khối: H = 1; C = 12; N = 14; O = 16; Na = 23; Mg = 24; Al = 27; S = 32; Cl = 35,5; K = 39; Ca = 40; Mn = 55; Fe = 56; Cu = 64; Zn = 65; Ag = 108; Cd = 112; Au = 197; Pb = 207; hằng số Faraday F = 96500 C/mol).}

\vspace{2pt}
\noindent
\textbf{PHẦN I (4,5 điểm -- 18 câu): Câu trắc nghiệm nhiều phương án lựa chọn.}\\
\textit{Thí sinh trả lời từ câu 1 đến câu 18. Mỗi câu hỏi thí sinh chỉ chọn một phương án.}

\setcounter{ex}{0}

\begin{ex}
(SBT – KNTT). Trong pin điện hoá Zn -- Cu, ở anode (cực âm) xảy ra quá trình
\choice
{oxi hoá Zn thành ion \ce{Zn^2+}}
{khử ion \ce{Cu^2+} thành Cu}
{khử Cu thành ion \ce{Cu^2+}}
{oxi hoá ion \ce{Zn^2+} thành Zn}
\loigiai{}
\end{ex}

\begin{ex}
Pin điện là gì?
\choice
{Một thiết bị lưu trữ và cung cấp điện năng}
{Một thiết bị dùng để đo nhiệt độ}
{Một thiết bị dùng để đo áp suất}
{Một thiết bị dùng để lưu trữ năng lượng cơ học}
\loigiai{}
\end{ex}

\begin{ex}
Pin Galvani còn được gọi là gì?
\choice
{Pin nhiên liệu}
{Pin điện hóa}
{Pin nhiệt điện}
{Pin quang điện}
\loigiai{}
\end{ex}

\begin{ex}
Điện cực dương trong pin gọi là gì?
\choice
{Anode}
{Cathode}
{Điện cực trơ}
{Cực âm}
\loigiai{}
\end{ex}

\begin{ex}
(SBT – KNTT). Trong quá trình mạ bạc cho một chiếc vòng bằng thép thì ở anode xảy ra quá trình
\choice
{\ce{Ag -> Ag+ + 1e}}
{\ce{Fe -> Fe^2+ + 2e}}
{\ce{2H2O -> 4H+ + O2 + 4e}}
{\ce{C -> C^4+ + 4e}}
\loigiai{}
\end{ex}

\begin{ex}
(SBT – KNTT). Khi điện phân dung dịch \ce{NaCl} bằng dòng điện một chiều (với điện cực trơ, có màng ngăn xốp) thì ở cathode xảy ra quá trình
\choice
{oxi hoá \ce{H2O} thành \ce{H+} và \ce{O2}}
{khử \ce{Cl-} thành \ce{Cl2}}
{oxi hoá \ce{Cl-} thành \ce{Cl2}}
{khử \ce{H2O} thành \ce{H2} và \ce{OH-}}
\loigiai{}
\end{ex}

\begin{ex}
(SBT – KNTT). Khi điện phân dung dịch \ce{CuSO4} bằng dòng điện một chiều (với điện cực anode bằng Cu) thì ở anode xảy ra quá trình
\choice
{oxi hoá \ce{H2O} thành \ce{H+} và \ce{O2}}
{khử \ce{Cu^2+} thành Cu}
{oxi hoá Cu thành \ce{Cu^2+}}
{khử \ce{H2O} thành \ce{H2} và \ce{OH-}}
\loigiai{}
\end{ex}

\begin{ex}
So sánh hai kim loại A và B có thế điện cực chuẩn lần lượt là $-0{,}25\text{ V}$ và $+0{,}34\text{ V}$, kim loại nào có khả năng khử mạnh hơn?
\choice
{Kim loại A}
{Kim loại B}
{Cả hai đều có khả năng khử như nhau}
{Không thể xác định}
\loigiai{}
\end{ex}

\begin{ex}
(SBT – CTST). Cho pin điện hoá Zn--Cu. Quá trình xảy ra ở cực dương của pin là
\choice
{\ce{Zn^2+ + 2e -> Zn}}
{\ce{Cu^2+ + 2e -> Cu}}
{\ce{Zn -> Zn^2+ + 2e}}
{\ce{Cu -> Cu^2+ + 2e}}
\loigiai{}
\end{ex}

\begin{ex}
(SBT – CTST). Cho pin điện hoá Al--Pb. Quá trình xảy ra ở cực âm của pin là
\choice
{\ce{Al^3+ + 3e -> Al}}
{\ce{Pb^2+ + 2e -> Pb}}
{\ce{Al -> Al^3+ + 3e}}
{\ce{Pb -> Pb^2+ + 2e}}
\loigiai{}
\end{ex}

\newpage

\begin{ex}
(SBT – CTST). Phát biểu nào sau đây \textbf{không} đúng về pin Galvani?
\choice
{Ở điện cực dương xảy ra quá trình oxi hoá}
{Cathode là kim loại yếu hơn, đóng vai trò điện cực dương của pin}
{Anode là kim loại mạnh hơn, đóng vai trò điện cực âm của pin}
{Phản ứng hóa học diễn ra trong pin kèm theo sự giải phóng điện năng}
\loigiai{}
\end{ex}

\begin{ex}
(SBT – CTST). Khi điện phân dung dịch nào sau đây, tại anode xảy ra quá trình oxi hóa nước?
\choice
{Dung dịch \ce{ZnCl2}}
{Dung dịch \ce{CuCl2}}
{Dung dịch \ce{AgNO3}}
{Dung dịch \ce{MgCl2}}
\loigiai{}
\end{ex}

\begin{ex}
(SBT – CTST). Điện phân dung dịch \ce{CuCl2} với điện cực trơ. Nồng độ của \ce{CuCl2} trong dung dịch
\choice
{giảm dần}
{tăng dần}
{không thay đổi}
{tăng dần rồi giảm dần}
\loigiai{}
\end{ex}

\begin{ex}
(SBT – CTST). Điện phân dung dịch \ce{CuSO4} với hai điện cực trơ. Sau một thời gian, màu xanh của dung dịch nhạt dần do:
\choice
{khí \ce{H2} sinh ra đã khử màu của dung dịch}
{dung dịch bị pha loãng dần nên màu xanh nhạt dần}
{ion \ce{Cu^2+} bị khử dần thành đồng kim loại}
{ion \ce{Cu^2+} được tạo thêm}
\loigiai{}
\end{ex}

\begin{ex}
(SBT – CTST). Điện phân điện cực trơ, có màng ngăn một dung dịch chứa các ion: \ce{Fe^2+}, \ce{Fe^3+}, \ce{Cu^2+}, \ce{Cl-}. Thứ tự điện phân xảy ra ở cathode là
\choice
{\ce{Fe^2+}, \ce{Fe^3+}, \ce{Cu^2+}}
{\ce{Fe^2+}, \ce{Cu^2+}, \ce{Fe^3+}}
{\ce{Fe^3+}, \ce{Cu^2+}, \ce{Fe^2+}}
{\ce{Fe^3+}, \ce{Fe^2+}, \ce{Cu^2+}}
\loigiai{}
\end{ex}

\begin{ex}
(SBT – CTST). Cho pin điện hoá Mn--Cd có $E^\circ_{\text{pin(Mn-Cd)}} = 0{,}79\text{ V}$ và $E^\circ_{\ce{Cd^2+/Cd}} = -0{,}40\text{ V}$. Thế điện cực chuẩn của cặp oxi hoá khử \ce{Mn^2+/Mn} có giá trị là
\choice
{0,27 V}
{-1,19 V}
{0,26 V}
{-0,25 V}
\loigiai{}
\end{ex}

\begin{ex}
(SBT – CTST). Quá trình điện phân dung dịch \ce{NaCl} (điện cực trơ, có màng ngăn) và điện phân \ce{NaCl} nóng chảy (điện cực trơ) có điểm giống nhau là
\choice
{tại cathode xảy ra sự khử ion \ce{Na+}}
{tại cathode xảy ra sự khử phân tử \ce{H2O}}
{tại anode xảy ra sự oxi hóa ion \ce{Cl-}}
{tại anode xảy ra sự oxi hóa phân tử \ce{H2O}}
\loigiai{}
\end{ex}

\begin{ex}
Cho sơ đồ phản ứng:\\
(1) $\ce{NaCl + H2O ->[\text{đpdd, cmn}] X + Y ^}\text{ (ở cathode)} + \ce{Z ^}\text{ (ở anode)}$\\
(2) $\ce{X + CO2 \text{ (dư)} -> T}$\\
Phát biểu nào sau đây đúng?
\choice
{Chất khí Z có thể khử được CaO thành Ca ở nhiệt độ cao}
{Dung dịch X có tính tẩy màu, sát trùng, thường gọi là Javel}
{Chất T được dùng làm thuốc giảm đau dạ dày}
{Chất khí Y không có màu, mùi, vị và Y có thể duy trì sự cháy}
\loigiai{}
\end{ex}

\newpage

\noindent
\textbf{PHẦN II (4,0 điểm -- 4 câu): Câu trắc nghiệm đúng sai.}\\
\textit{Thí sinh trả lời từ câu 1 đến câu 4. Trong mỗi ý a), b), c), d) ở mỗi câu, thí sinh chọn đúng hoặc sai.}

\setcounter{ex}{0}

\begin{ex}
(SBT – CĐ). Chọn tính đúng sai cho mỗi phát biểu sau về acquy:
\choiceTF
{Một ưu điểm của acquy là tái sử dụng được nhiều lần}
{Phản ứng xảy ra trong acquy cũng giống như phản ứng xảy ra trong pin Galvani nhưng có thể đảo ngược}
{Acquy không gây ô nhiễm môi trường}
{Acquy là nguồn điện hoá học có thể hoạt động liên tục}
\loigiai{}
\end{ex}

\begin{ex}
(SBT – CTST). Em hãy cho biết những thí nghiệm sau đây có thể tạo thành một pin điện hóa là đúng hay sai:
\choiceTF
{Một điện cực nhôm và một điện cực đồng cùng nhúng vào dung dịch nước muối}
{Một điện cực nhôm và một điện cực đồng cùng nhúng vào nước cất}
{Hai điện cực làm bằng đồng cùng nhúng vào dung dịch nước vôi trong}
{Hai điện cực làm bằng nhựa nhúng vào dầu hoả}
\loigiai{}
\end{ex}

\begin{ex}
(SBT – KNTT). Trong quá trình điện phân dung dịch \ce{CuSO4} với anode bằng đồng. Chọn tính đúng sai cho mỗi phát biểu sau:
\choiceTF
{Ở anode xảy ra quá trình oxi hóa nước}
{Khối lượng anode không thay đổi}
{Nồng độ \ce{CuSO4} trong dung dịch giảm dần}
{Khối lượng cathode tăng}
\loigiai{}
\end{ex}

\begin{ex}
(SBT – KNTT). Điện phân dung dịch \ce{NaCl} bão hoà (với điện cực trơ, màng ngăn xốp) đến khi nồng độ \ce{NaCl} giảm đi một nửa thì dừng điện phân. Chọn tính đúng sai cho mỗi phát biểu sau:
\choiceTF
{Dung dịch sau điện phân làm phenolphthalein chuyển màu hồng}
{Ở cathode chỉ xảy ra quá trình khử ion \ce{Na+}}
{Số mol khí \ce{Cl2} thoát ra ở anode bằng số mol \ce{H2} thoát ra ở cathode}
{Thứ tự điện phân ở anode là \ce{H2O}, \ce{Cl-}}
\loigiai{}
\end{ex}

\newpage

\noindent
\textbf{PHẦN III (1,5 điểm -- 6 câu): Câu trắc nghiệm yêu cầu trả lời ngắn.}\\
\textit{Thí sinh trả lời từ câu 1 đến câu 6.}

\setcounter{ex}{0}

\begin{ex}
Cho bảng số liệu thế điện cực chuẩn: \ce{Al^3+/Al} ($-1{,}676\text{ V}$); \ce{Zn^2+/Zn} ($-0{,}763\text{ V}$); \ce{Fe^2+/Fe} ($-0{,}440\text{ V}$); \ce{Cu^2+/Cu} ($+0{,}340\text{ V}$); \ce{Ag+/Ag} ($+0{,}799\text{ V}$); \ce{Au^3+/Au} ($+1{,}520\text{ V}$). Hãy cho biết có bao nhiêu kim loại trong bảng có tính khử mạnh hơn Cu?
\par\noindent\textit{Trả lời:} \dotfill
\loigiai{}
\end{ex}

\begin{ex}
Cho pin điện hóa tạo bởi hai cặp oxi hoá -- khử \ce{Zn^2+/Zn} và \ce{Cu^2+/Cu}. Cho $E^\circ_{\ce{Cu^2+/Cu}} = +0{,}34\text{ V}$ và $E^\circ_{\ce{Zn^2+/Zn}} = -0{,}76\text{ V}$. Tính sức điện động chuẩn của pin (theo đơn vị V).
\par\noindent\textit{Trả lời:} \dotfill
\loigiai{}
\end{ex}

\begin{ex}
Cho các kim loại sau: Na, Ca, Fe, Al, Zn, Cu, Ag. Số kim loại được điều chế trong công nghiệp bằng phương pháp điện phân hợp chất nóng chảy của chúng là bao nhiêu?
\par\noindent\textit{Trả lời:} \dotfill
\loigiai{}
\end{ex}

\begin{ex}
Cho các chất rắn sau đây: \ce{NaCl}, \ce{CaCO3}, \ce{CuSO4}, \ce{NaOH}, \ce{CaCl2}, \ce{FeCl2}, \ce{CuCl2}, \ce{MgCl2}, \ce{Al2O3}. Số chất có thể dùng để điều chế kim loại bằng phương pháp điện phân nóng chảy là bao nhiêu?
\par\noindent\textit{Trả lời:} \dotfill
\loigiai{}
\end{ex}

\begin{ex}
Một viên pin lithium-ion của điện thoại có điện áp danh định 3,7 V và dung lượng 4000 mAh. Điện thoại tiêu thụ trung bình dòng điện 500 mA khi sử dụng liên tục. Hỏi điện thoại có thể hoạt động liên tục trong bao lâu trước khi cần sạc lại (theo đơn vị giờ)?
\par\noindent\textit{Trả lời:} \dotfill
\loigiai{}
\end{ex}

\begin{ex}
Một nhà máy muốn mạ một lớp đồng lên một tấm kim loại có diện tích $1\text{ m}^2$ để chống ăn mòn. Họ sử dụng dung dịch \ce{CuSO4} và dòng điện có cường độ 500 A trong thời gian 2 giờ. Biết khối lượng mol của Cu là 64 g/mol. Hãy tính khối lượng đồng (theo đơn vị kg) bám vào tấm kim loại. (Làm tròn kết quả đến hàng phần trăm).
\par\noindent\textit{Trả lời:} \dotfill
\loigiai{}
\end{ex}


\Closesolutionfile{ans}

\end{document}
